% tau_field.tex -- round 393, lemma-first attack on
% Delta^2 log tau under F1 (the F_eps rest named by r392).
% Outcome: REDUCED.  Builds with pdflatex.
% Companion machine check: verify_tau_field.py.
% Discovery census: experiments/tfpt-discovery/tau_field_probe.py.
% NO RH CLAIM.  Research documentation, not a theorem of RH.
\documentclass[11pt]{article}

\usepackage[a4paper,margin=2.7cm]{geometry}
\usepackage{amsmath,amssymb,amsthm}
\usepackage{booktabs}
\usepackage{microtype}
\usepackage{xcolor}
\usepackage[colorlinks=true,linkcolor=blue!50!black,urlcolor=blue!50!black]{hyperref}

\newtheorem{lemma}{Lemma}
\newtheorem{prop}{Proposition}
\newtheorem{definition}{Definition}
\theoremstyle{remark}
\newtheorem{remark}{Remark}

\newcommand{\N}{\mathbb{N}}
\newcommand{\Q}{\mathbb{Q}}
\newcommand{\R}{\mathbb{R}}
\newcommand{\eps}{\varepsilon}

\title{\bfseries $\Delta^2\log\tau$ under $F_1$, reduced\\[2mm]
\large Cluster anatomy of the Uvarov $\tau$-field,
rank-one locality, and why $F_1$ does not close $F_\eps$}
\author{Stefan Hamann}
\date{August 28, 2026}

\begin{document}
\maketitle

\begin{abstract}\noindent
Round 392 reduced $F_\eps$ to a bound on $\Delta^2\log\tau$ for
$F_1$-bounded $\nu$-runs (max run $\le 2$, r382, MAIN $85/85$),
with the box $\lvert e^{\Delta^2\log\tau}-1\rvert\le\tfrac14$
equivalent to $\lvert\gamma-\tfrac14\rvert\le\tfrac1{16}$ on
discrete Bernstein--Szeg\H{o}, and the last-12 jump of $\gamma$
equal to $\lvert\Delta(\Delta^2\log\tau)\rvert$.
This note splits $\Xi$ under $F_1$ into singletons and pairs,
writes $\tau$ as an explicit $1\times1$/$2\times2$ product times
a coupling determinant $\kappa$, and proves that $\Delta^2\log\tau$
is the first difference of a rank-one increment
(matrix-determinant lemma) --- the cluster-count does not enter
linearly.
What is proved: the decomposition identity; rank-one locality;
coupling is load-bearing ($\sim 31\%$ of $\lvert\Delta^2\rvert$ on
FRAME-A $w=9$); a consecutive run of length $3$ is a $3\times3$
block that leaves the pair bound (jump $0.4619$, out even at the
legal $V_3'$ relaxation $\mathrm{JUMP}'=0.45$);
two-period AP has $\Delta^2=0$.
What is \emph{not} proved: the bound itself.
$F_1$ is necessary and not sufficient --- a random $F_1$
half-filling on the cosine grid leaves every legal jump
constant (h$=80$, jump $1.18$); the cluster $L^1$ triangle
($0.436>\log\tfrac54$) does not close the box; raising JUMP to
$0.45$ covers the isolated-pair toy ($0.4183$) and the kz~$55$
census row ($0.3942$) and does not make either a theorem of $F_1$.
No RH claim.
\end{abstract}

\medskip\noindent
\textbf{Status.}\enspace
Lemmas~\ref{lem:sm}--\ref{lem:rank1} are proved (finite identities).
Lemma~\ref{lem:w9} is a machine identity.
Propositions~\ref{prop:f1}--\ref{prop:red} name the reduction.
No RH claim.  No $L^*$ claim.  No $G_\varepsilon$ claim.

\section{Recollection}\label{sec:rec}

Occupied Fej\'er is node deletion of the $\nu$-support $\Xi$ from
the discrete Bernstein--Szeg\H{o} cosine grid~\cite{r392,r390}.
Writing $\tau_n=\det(I-K_n[\Xi]W)$ with $\tau_0=1$,
\[
\gamma_n(\mu)
=
\gamma_n(\sigma)
\cdot
\frac{\tau_{n+1}\,\tau_{n-1}}{\tau_n^2}.
\]
On $\gamma_n(\sigma)=\tfrac14$ the box
$\lvert\gamma-\tfrac14\rvert\le\tfrac1{16}$ is
$\lvert e^{\Delta^2\log\tau}-1\rvert\le\tfrac14$, and the last-12
jump of $\gamma$ is $\lvert\Delta(\Delta^2\log\tau)\rvert$.
$F_1$ (r382) is $\max$ consecutive $\nu$-run $\le 2$ on the cosine
grid, MAIN $85/85$.
$V_3'$/$A_{15}$ first fails at coherent last-12 scale $2.0$;
scale $1.5$ stays inside~\cite{v3p}.
Raising $\mathrm{JUMP}$ from $\tfrac25$ to $0.45$ is therefore
legal air; $\log\tfrac53\approx 0.511$ would make the jump
redundant with the box and is not used.

\section{Cluster anatomy}\label{sec:cluster}

Under $F_1$, $\Xi$ is a disjoint union of singletons and adjacent
pairs.  Write $K_n$ for the CD Gram of $\sigma$ on $\Xi$.

\begin{lemma}[$1\times1$ and $2\times2$]\label{lem:sm}
A singleton $\xi$ of mass $w$ contributes
\[
\tau_n^{(\xi)}=1-w\,K_n(\xi,\xi)
\]
(Sherman--Morrison).
An adjacent pair $(\xi_1,\xi_2)$ of masses $(w_1,w_2)$ contributes
the $2\times2$ determinant
\begin{align*}
\tau_n^{(\xi_1\xi_2)}
&=
(1-w_1 K_{11})(1-w_2 K_{22})-K_{12}^2 w_1 w_2.
\end{align*}
Exact over $\Q$ on a five-atom toy (singleton $1/5$; adjacent pair
$2/15$).
\end{lemma}

\begin{lemma}[decomposition]\label{lem:decomp}
Writing $\tau_n^{\mathrm{diag}}=\prod_c\tau_n^{(c)}$ over the
$F_1$ clusters and $\kappa_n=\tau_n/\tau_n^{\mathrm{diag}}$,
\[
\Delta^2\log\tau
=
\sum_c\Delta^2\log\tau^{(c)}
+
\Delta^2\log\kappa.
\]
The identity holds at $3\cdot10^{-14}$ on FRAME-A $w=9$
(last $20$ degrees).
Two far singletons give $\kappa=5/9$; two neighbours give a
strictly larger $\lvert\log\kappa\rvert$ (CD decay of $K_{12}$).
\end{lemma}

On $w=9$: $n_\nu=104$ nodes split as $102$ clusters
($100$ singletons, $2$ pairs), $F_1$ true.
Last-12 $\max\lvert\Delta^2\log\tau\rvert=0.1553$,
$\lvert e^{\Delta^2}-1\rvert_{\max}=0.1438<\tfrac14$,
$\lvert\Delta(\Delta^2)\rvert_{\max}=0.2379$ equal to the
occupied-Fej\'er jump.
The diagonal piece is $0.1462$; the coupling piece is $0.0490$
($\sim 31\%$ of the full last-12).
Coupling is load-bearing: omitting $\Delta^2\log\kappa$ is a named
kill (error $0.049$).

\section{Rank-one locality}\label{sec:rank1}

The cluster count scales with $N$.  A triangle bound
$\sum_c\lvert\Delta^2\log\tau^{(c)}\rvert$ is therefore the wrong
object unless $\Delta^2$ cancels the volume.

\begin{lemma}[rank-one increment]\label{lem:rank1}
$K_{n+1}=K_n+\pi_n\pi_n^\top$.  The matrix-determinant lemma gives
\[
\rho_n
:=
\frac{\tau_{n+1}}{\tau_n}
=
1-(W\pi_n)^\top(I-K_n W)^{-1}\pi_n,
\]
hence
\[
\Delta^2\log\tau_n
=
\log\rho_{n+1}-\log\rho_n.
\]
Exact at $9\cdot10^{-15}$ on $w=9$.
Adjacent degrees share every cluster contribution except the
rank-one update at $n$: $\Delta^2$ reduces the volume sum to a
local increment in the degree window.  That is the structural
statement of this round.  It does \emph{not} by itself bound
$\lvert\log\rho_{n+1}/\rho_n\rvert$.
\end{lemma}

On $w=9$ the last-12 $L^1$ of the cluster pieces is $0.436$,
against $\lvert\sum\rvert=0.066$ (cancellation $0.85$) and against
$\log\tfrac54\approx 0.223$.
The triangle does not close the box.
Per-cluster last-12 maxabs has median $0.0072$ and max $0.0125$;
$n_{\mathrm{cl}}\times\mathrm{med}\approx 0.73$ is five times the
coherent sum.

\section{Adversaries}\label{sec:adv}

\begin{lemma}[$w=9$ pins]\label{lem:w9}
FRAME-A $w=9$ satisfies $F_1$ with the cluster split of
Lemma~\ref{lem:decomp}; last-12 occupied Fej\'er lies in the box
at jump $0.2379$.  Machine identity against the Uvarov formula
at $9\cdot10^{-15}$ is r392.
\end{lemma}

\begin{prop}[$F_1$ necessary, not sufficient]\label{prop:f1}
A consecutive $\nu$-run of length $3$ on $\mathtt{full\_grid}(40)$,
last $12$ of $n\le 60$, has jump $0.4619$ and
$\lvert\gamma-\tfrac14\rvert=0.0672>\tfrac1{16}$: out at
$\mathrm{JUMP}=\tfrac25$ and at $\mathrm{JUMP}'=0.45$.
The $3\times3$ cluster $\Delta^2$ on a height-$12$ toy exceeds
the pair $2\times2$ ($\lvert e^{\Delta^2}-1\rvert=0.630>0.25$).
$F_1$ is therefore necessary for the bound.
It is not sufficient.
Periodic $F_1$ tilings ($1100$, $1010$, $110$) stay in at the
same window; a random $F_1$ half-filling on $\mathtt{full\_grid}(80)$
at $n=60$ has jump $1.179$ and leaves $\mathrm{JUMP}'$.
The isolated central pair (F1-legal, not half-filled) has jump
$0.4183$: out at $\tfrac25$, in at $0.45$, inside the $1/16$ box.
Scramble occupation seed $3$ has $\max$ run $6$ ($12$ runs of
length $\ge 3$), $F_1$ false, occupied jump $0.4438$: a cluster
break, not a coupling remainder.
Core-$42$ occupied last-12: $0/42$ out at $\tfrac25$, max jump
$0.3942$ at kz~$55$ --- census of the $L^*$-folded class, not a
corollary of $F_1$.
\end{prop}

Two-period AP deletion makes every cluster contribution identical
across $n$, so $\Delta^2\log\tau=0$, consistent with the r392
non-kill (still Bernstein--Szeg\H{o}).
Wrong Sherman--Morrison sign on the cluster formula is a mutant
(last-12 error $0.199$).

\section{What remains}\label{sec:rest}

\begin{prop}[reduction]\label{prop:red}
The $\tau$-field under $F_1$ is SATZ as algebra
(Lemmas~\ref{lem:sm}--\ref{lem:rank1}).
$F_\eps$ does \emph{not} follow from $F_1$ + half-filling + mesh
separation: random $F_1$ half-fillings leave every legal jump
constant, and the cluster triangle does not close.
Raising JUMP to $0.45$ (legal $V_3'$ air) covers the isolated-pair
toy and the kz~$55$ margin $0.0058$, and still leaves $F_\eps$ a
census of the $L^*$-folded occupation, not a theorem of $F_1$.
The remaining named rest is $L^*$-occupation regularity --- the
gap between $F_1$-legal cosine masks and the folded half-filling
that actually occurs --- sibling to the Sign-Schur lemma of~\cite{r392}
for Assist.
The chain
\[
F_\eps
\;\Longleftarrow\;
\bigl(\Delta^2\log\tau\text{ under }F_1\bigr)
\;\Longleftarrow\;
\bigl(L^*\text{-regular }F_1\text{ occupation}\bigr)
\]
is the honest reduction.
$W_\eps$ is untouched.
No $m_4$.  No RH claim.
\end{prop}

\section{Machine checks}\label{sec:machine}

Numbered lemmas: \texttt{verify\_tau\_field.py} (same directory).
Standalone: $1\times1$/$2\times2$ over $\Q$; $\kappa$ vs gap;
$\tau_0=1$; JUMP air.
Construction: $w=9$ $F_1$ split, rank-one, decomposition, coupling,
AP non-kill, run-$3$ kill, scramble $F_1$-break, coupling mutant.
Discovery census: \texttt{tau\_field\_probe.py}
(\texttt{--smoke} / full), core-$42$, random-$F_1$ counterexample.
Exit line: \texttt{TAU FIELD VERIFIED}.

\section*{Claim boundary}

Finite identities for the Uvarov $\tau$-field of a positive atomic
measure on a cosine grid under $F_1$-bounded runs, a named
reduction of $F_\eps$ to $L^*$-occupation regularity, and named
clustered / random-$F_1$ / scramble breaks.
No statement about zeros of $\zeta(s)$, in either direction.
No RH claim.

\begin{thebibliography}{9}
\bibitem{r392}
S.~Hamann,
\textit{The deletion transform, reduced},
standalone note, August 2026,
\texttt{rh/problem/deletion\_transform.tex}.
\bibitem{r390}
S.~Hamann,
\textit{The $G_\varepsilon^\mu$-lemma, reduced},
standalone note, August 2026,
\texttt{rh/problem/g\_eps\_mu.tex}.
\bibitem{r382}
S.~Hamann,
\textit{The pivot-entry lemma},
standalone note, August 2026,
\texttt{rh/problem/pivot\_entry\_lemma.tex}.
\bibitem{v3p}
S.~Hamann,
\textit{The $V_3'$-lemma, reduced},
standalone note, August 2026,
\texttt{rh/problem/v3prime\_proof.tex}.
\end{thebibliography}

\end{document}
